% AVM Regional Reliability Report
% Compile with pdflatex (portable): pdflatex AVM_Report.tex
% Province names are romanised (ASCII) in this document so it compiles with a
% plain TeX Live install; the per-region CSV files keep the original Vietnamese names.
% This file \input{}s gen_region_metrics_accuracy.tex, gen_region_metrics_calibration.tex
% and gen_region_examples.tex, all generated by src/build_regional_report.py.
\documentclass[11pt]{article}

\usepackage[utf8]{inputenc}
\usepackage[T1]{fontenc}
\usepackage[margin=0.85in]{geometry}
\usepackage{booktabs}
\usepackage{array}
\usepackage{longtable}
\usepackage{amsmath}
\usepackage{xcolor}
\usepackage{hyperref}
\hypersetup{colorlinks=true, linkcolor=blue!55!black, urlcolor=blue!55!black}

\newcommand{\pct}{\%}
\definecolor{good}{RGB}{20,120,70}
\definecolor{bad}{RGB}{170,40,40}

\title{\textbf{Automated Valuation Model for Vietnamese Landed Residential Houses}\\[4pt]
\large A Region-Segmented Reliability Report and Production Defense}
\author{AVM Project}
\date{\today}

\begin{document}
\maketitle

\begin{abstract}
\noindent This report documents an Automated Valuation Model (AVM) for Vietnamese landed
residential houses, from raw-data preprocessing through the deployed model and its results.
Its central argument is that a single dataset-wide accuracy number is the wrong lens for a real
product: accuracy must be read \emph{per region}, large errors must be \emph{explained} rather than
only minimised, and the system must \emph{know when it cannot answer}. We show that the typical
listing is valued to about $15\pct$ error, that the model is statistically unbiased, that its large
errors concentrate in a thin population of mislabeled/anomalous listings (not ordinary homes), and
that a confidence-tiered policy lets the product serve normal listings sharply while flagging
anomalies. We close by diving into each region and, for the region that produces the most errors
(Ho Chi Minh City), contrasting its worst error case against a normal case.
\end{abstract}

\tableofcontents
\bigskip
\hrule
\bigskip

\section*{Key Insights at a Glance}
\begin{enumerate}
  \item \textbf{Typical accuracy is strong.} The production model values the typical listing to
        \textbf{$\sim 15\pct$} error, with \textbf{$\sim 59\pct$} of estimates within $20\pct$.
  \item \textbf{The model is unbiased.} The alarming raw bias ($+13.9\pct$) is a mean-of-ratios
        artifact; the median bias is $\approx 0$ overall \emph{and in every region}.
  \item \textbf{Big errors are bad labels, not bad predictions.} Large errors concentrate in listings
        with an implausibly low price/m\textsuperscript{2} (mistyped / underpriced posts); the model
        actually predicts the sensible value and thus \emph{catches} them.
  \item \textbf{The model knows when it is unsure.} Its own interval width predicts its error
        ($8\pct$ median APE on the most-confident fifth vs.\ $29\pct$ on the least), enabling an
        auto-serve / review / flag policy.
  \item \textbf{Reliability is reported per region.} Da Nang, Khanh Hoa, Hai Phong are product-ready
        ($\sim 11$--$13\pct$); \textbf{Ho Chi Minh} is the weak spot ($23.7\pct$) and the clearest
        next target.
\end{enumerate}
\bigskip
\hrule

%======================================================================
\section{Data Preprocessing}
%======================================================================
Raw listings are scraped, heterogeneous, and noisy. The following steps turn them into a clean
modeling table (implemented in \texttt{src/common.py} and the pipeline scripts):

\begin{enumerate}
  \item \textbf{Parsing and typing.} Extract structured numeric fields (area in m\textsuperscript{2},
        number of floors, frontage, road width, bedrooms/bathrooms, width, length, built year) and
        parse dates (post date, updated date) into a listing year.
  \item \textbf{Address normalization.} Merge province/district spelling variants and abbreviations
        into a single canonical label --- e.g.\ \texttt{tp.hcm}, \texttt{tphcm}, and
        \texttt{ho chi minh} collapse to one region; tone variants (\texttt{hoa}$\rightarrow$\texttt{hoa},
        etc.) are unified. This stops one real place from splitting into several location buckets.
  \item \textbf{Property classification.} Assign each listing to the landed-residential family and
        one of seven subtypes (frontage, alley, normal, alley-frontage, townhouse-project, villa,
        shophouse) from text and structured signals.
  \item \textbf{Text feature extraction.} Derive boolean flags from the free-text description
        (frontage, alley, business, car access, corner, villa, legal, project, urban area) plus text
        length and word count.
  \item \textbf{Near-duplicate removal.} Collapse the same listing re-posted many times ---
        matched on address, an area bucket ($\pm 3$ m\textsuperscript{2}), and a price bucket
        ($\pm 200$M VND) --- keeping only the most recent copy, so repeated posts do not dominate
        training.
  \item \textbf{Suspicious-price auditing.} Flag and remove implausible prices (placeholder
        sequences, sequential/repeated digit runs, sub-500M values whose text says ``ty''/billions),
        removing roughly $8{,}000$ rows in the cleaned variant.
  \item \textbf{Recency weighting.} Assign a training weight by listing year ($2025\!\rightarrow\!1.0$,
        older years progressively lower, undated $\rightarrow 0.5$) to down-weight stale prices.
  \item \textbf{Leakage-safe outlier bands.} Compute local price-per-m\textsuperscript{2} bands on
        the \emph{training} data only, used to weight/robustify fitting without leaking test
        information.
\end{enumerate}

%======================================================================
\section{Data Used to Train the Model}
%======================================================================
The cleaned data is split by a stratified random split on subtype:

\begin{table}[h]
\centering
\begin{tabular}{lr}
\toprule
\textbf{Stage} & \textbf{Rows} \\
\midrule
Cleaned listings & 813,215 \\
Modeling rows (landed residential, valid price/area) & 704,370 \\
Training set & 563,496 \\
Test set (held out) & 140,874 \\
\bottomrule
\end{tabular}
\caption{Data volumes through the pipeline.}
\end{table}

\noindent\textbf{Subtype mix (modeling rows):} frontage 171k, normal 152k, alley 80k,
alley-frontage 55k, townhouse-project 52k, villa 44k, shophouse 10k. The target is
$\log(\text{price})$; models are gradient-boosted trees (LightGBM, L1 objective) and predictions are
back-transformed with $\exp(\cdot)$.

%======================================================================
\section{Results}
%======================================================================
All models are evaluated on the same 140,874-row held-out test set.

\begin{table}[h]
\centering
\begin{tabular}{lccc l}
\toprule
\textbf{Model} & \textbf{Median APE} & \textbf{Hit\,$\leq$20\pct} & \textbf{MAPE} & \textbf{Role} \\
\midrule
Global               & 21.96\pct & 46.5\pct & 44.0\pct & baseline \\
Subtype              & 20.45\pct & 49.2\pct & 42.6\pct & one model / subtype \\
Region               & 15.36\pct & 59.0\pct & 34.1\pct & one model / province \\
Hybrid (hard router) & 15.39\pct & 58.3\pct & 36.1\pct & previous production \\
\textbf{Shrinkage blend} & \textbf{14.84\pct} & \textbf{59.5\pct} & \textbf{34.5\pct} & \textbf{current production} \\
Region$\times$Subtype & 14.01\pct & 61.3\pct & 32.1\pct & best raw (dense groups) \\
\bottomrule
\end{tabular}
\caption{Model comparison. \emph{Median APE} is the typical listing's error and is the headline
metric for a skewed target; \emph{MAPE} (a mean) is inflated by a heavy right tail.}
\end{table}

\paragraph{The production model.} The \textbf{shrinkage blend} is the deployed model: it values
\emph{every} listing (unlike Region$\times$Subtype, which only covers dense groups) and it improves on
the previous hard router by $-0.55$ Median-APE points and $+1.14$ hit-rate points --- a statistically
solid gain (McNemar $p<0.001$; net $+1{,}607$ listings moved inside the $\pm 20\pct$ band).

\paragraph{The model is unbiased --- read this before judging the numbers.}
The raw mean bias looks large ($+13.9\pct$), but that is a \emph{reporting artifact}. Because the
model predicts the median, $\text{median}(\text{pred}/\text{true})\approx 1.00$; the robust bias
views confirm calibration:

\begin{table}[h]
\centering
\begin{tabular}{lc}
\toprule
\textbf{Bias measure (Region$\times$Subtype)} & \textbf{Value} \\
\midrule
Raw mean bias & $+13.9\pct$ \\
Median bias & $-0.29\pct$ \\
10\pct-trimmed mean bias & $+0.63\pct$ \\
\bottomrule
\end{tabular}
\caption{The model is centered; the raw mean is pulled up by a thin right tail (Section 6).
Forcing the raw bias to zero would push Median APE from $12.8\pct$ to $16.8\pct$ --- so we
deliberately do not apply a naive bias correction.}
\end{table}

Each valuation also carries a \textbf{P10--P90 prediction interval}; empirical coverage on the test
set is $76.7\pct$ (target $80\pct$). This interval is the foundation of the production defense.

%======================================================================
\section{Interpretation for a Real Product}
%======================================================================
Additional model tuning now moves the headline error by only fractions of a point. For a real
product the useful question is not ``what is the average error?'' but:
\begin{quote}
\emph{For which listings, and in which markets, can we trust the valuation --- and what does the
system do when it cannot?}
\end{quote}
This reframes the deliverable around three business-relevant claims, each defended with evidence:
\begin{enumerate}
  \item \textbf{Reliability is segmented.} Accuracy is reported \emph{per region}, because a national
        average hides that some markets are prediction-ready today and others are data-limited
        (Section~\ref{sec:regional}).
  \item \textbf{Errors are explained, not just minimised.} Large errors concentrate in a thin,
        \emph{identifiable} population of anomalous / mislabeled listings, not ordinary homes
        (Section~\ref{sec:errors}).
  \item \textbf{Behaviour on new cases is defensible.} A normal incoming listing is valued sharply;
        the model \emph{knows when it cannot} (via its own interval) and flags anomalies rather than
        emitting a confident wrong number (Section~\ref{sec:defense}).
\end{enumerate}

\subsection{The model knows when it is uncertain}\label{sec:defense}
Grouping listings by the width of their own predicted interval --- a signal available \emph{before}
the true price is known --- accuracy falls monotonically:

\begin{table}[h]
\centering
\begin{tabular}{lccc}
\toprule
\textbf{Interval-width quintile} & \textbf{n} & \textbf{Median APE} & \textbf{Hit\,$\leq$20\pct} \\
\midrule
Q1 (narrowest) & 28,175 & 8.05\pct  & 81.1\pct \\
Q2             & 28,175 & 11.48\pct & 68.9\pct \\
Q3             & 28,174 & 15.47\pct & 59.0\pct \\
Q4             & 28,175 & 20.24\pct & 49.7\pct \\
Q5 (widest)    & 28,175 & 28.85\pct & 38.7\pct \\
\bottomrule
\end{tabular}
\caption{The model's self-reported uncertainty is a strong, honest predictor of its real error.}
\end{table}

This yields a three-tier operating policy (no true price required):

\begin{table}[h]
\centering
\begin{tabular}{llccc}
\toprule
\textbf{Tier} & \textbf{Action} & \textbf{Share} & \textbf{Median APE} & \textbf{Hit\,$\leq$20\pct} \\
\midrule
A --- auto-serve      & return value automatically      & 40\pct & 9.55\pct  & 75.0\pct \\
B --- serve with range& return value with P10--P90 band & 40\pct & 17.64\pct & 54.3\pct \\
C --- flag for review & route to human review           & 20\pct & 28.85\pct & 38.7\pct \\
\bottomrule
\end{tabular}
\caption{Sharp where confident, cautious where not. A complementary check: when a listing's asking
price falls \emph{inside} its interval ($76.7\pct$ of listings) Median APE is $10.9\pct$; when it
falls \emph{outside} it is $40.0\pct$.}
\end{table}

\noindent The dangerous failure mode of an AVM is not ``sometimes wrong'' --- every AVM is --- but
``\emph{confidently} wrong without warning.'' This system avoids that: its error is concentrated in a
population it can identify up front.

%======================================================================
\section{Diving Into Each Region}\label{sec:regional}
%======================================================================
Every province with $\geq 100$ test listings has its own CSV file named after it (e.g.\
\texttt{regions/ho chi minh.csv}) holding all its listings with actual price, prediction, interval,
percentage error, and a \texttt{case\_type} label, sorted worst-error first. The 48 reported
provinces cover $\sim 99\pct$ of listings.

All metrics are reported for every region, split across two tables for readability:
Table~\ref{tab:regacc} covers \textbf{accuracy} (sample size, MAE, MAPE, median APE, and the hit
rates within 10\pct\ and 20\pct), and Table~\ref{tab:regcal} covers \textbf{calibration and
confidence} (raw / median / trimmed bias, mean interval width, in-band coverage, and the
high-confidence median APE). Both are sorted by sample size.

{\scriptsize
\setlength{\tabcolsep}{4pt}
\begin{longtable}{lrrrrrr}
\caption{Accuracy metrics per region (production shrinkage-blend model).}\label{tab:regacc}\\
\toprule
\textbf{Region} & \textbf{n} & \textbf{MAE (B)} & \textbf{MAPE\%} & \textbf{Med.APE\%} & \textbf{Hit10\%} & \textbf{Hit20\%} \\
\midrule
\endfirsthead
\multicolumn{7}{c}{\tablename\ \thetable\ -- continued}\\
\toprule
\textbf{Region} & \textbf{n} & \textbf{MAE (B)} & \textbf{MAPE\%} & \textbf{Med.APE\%} & \textbf{Hit10\%} & \textbf{Hit20\%} \\
\midrule
\endhead
\midrule\multicolumn{7}{r}{\itshape continued on next page}\\
\endfoot
\bottomrule
\endlastfoot
Da Nang & 26,787 & 2.40 & 21.0 & 11.3 & 45.9 & 69.5 \\
Ho Chi Minh & 14,904 & 5.12 & 48.0 & 23.6 & 24.4 & 43.7 \\
Hai Phong & 13,952 & 1.57 & 23.8 & 12.8 & 42.3 & 65.1 \\
Binh Duong & 13,663 & 1.19 & 30.5 & 12.8 & 42.8 & 64.0 \\
Dong Nai & 12,978 & 1.70 & 32.0 & 14.6 & 39.5 & 59.2 \\
Khanh Hoa & 8,034 & 2.66 & 24.8 & 11.2 & 47.1 & 66.5 \\
Ba Ria Vung Tau & 6,873 & 2.79 & 27.9 & 11.9 & 45.5 & 64.6 \\
Long An & 6,496 & 1.07 & 33.5 & 17.7 & 33.4 & 53.7 \\
Can Tho & 3,514 & 2.40 & 30.0 & 13.3 & 42.3 & 62.1 \\
Ha Noi & 3,358 & 6.52 & 152.5 & 20.4 & 28.6 & 49.3 \\
Hue & 2,566 & 1.91 & 25.6 & 11.9 & 44.9 & 66.0 \\
Dak Lak & 2,228 & 1.45 & 32.1 & 14.0 & 39.6 & 62.2 \\
Lam Dong & 2,175 & 3.98 & 36.0 & 18.2 & 33.2 & 52.5 \\
Bac Ninh & 2,061 & 2.87 & 33.4 & 14.5 & 39.4 & 60.1 \\
Quang Nam & 1,896 & 3.55 & 32.0 & 17.6 & 33.0 & 54.4 \\
Quang Ninh & 1,565 & 5.38 & 45.8 & 18.0 & 33.1 & 53.2 \\
Binh Thuan & 1,518 & 3.12 & 39.2 & 20.0 & 31.7 & 49.9 \\
Tay Ninh & 1,166 & 1.20 & 35.7 & 22.1 & 29.8 & 46.4 \\
Thanh Hoa & 1,150 & 2.46 & 42.5 & 20.6 & 32.4 & 49.0 \\
other small (46) & 936 & 3.25 & 110.6 & 53.1 & 10.6 & 21.6 \\
Hoa Binh & 912 & 2.04 & 44.2 & 23.0 & 27.2 & 43.9 \\
Nghe An & 875 & 2.43 & 40.8 & 15.1 & 37.7 & 59.7 \\
Vinh Phuc & 843 & 2.38 & 39.4 & 18.0 & 33.9 & 54.0 \\
Binh Dinh & 836 & 1.98 & 28.4 & 14.4 & 38.8 & 60.4 \\
Nam Dinh & 815 & 1.11 & 23.0 & 13.7 & 39.6 & 63.7 \\
Thai Binh & 813 & 1.27 & 28.1 & 15.6 & 35.8 & 58.3 \\
Hai Duong & 712 & 2.60 & 29.6 & 13.8 & 40.9 & 60.5 \\
Binh Phuoc & 663 & 1.24 & 44.7 & 27.1 & 23.1 & 39.7 \\
Bac Giang & 582 & 1.90 & 44.8 & 26.4 & 22.3 & 41.1 \\
Tien Giang & 560 & 1.51 & 38.5 & 21.6 & 30.7 & 48.0 \\
Kien Giang & 531 & 5.02 & 55.7 & 26.9 & 25.8 & 42.6 \\
Hung Yen & 503 & 3.59 & 33.5 & 17.8 & 32.2 & 53.1 \\
Phu Yen & 454 & 1.65 & 36.0 & 19.6 & 29.3 & 50.4 \\
Ha Nam & 448 & 2.17 & 38.1 & 17.7 & 33.0 & 52.7 \\
Tra Vinh & 331 & 0.87 & 34.9 & 22.5 & 22.7 & 45.6 \\
Ninh Thuan & 323 & 1.53 & 55.3 & 20.5 & 32.5 & 49.2 \\
Ben Tre & 318 & 1.94 & 42.3 & 22.7 & 29.6 & 45.6 \\
Thai Nguyen & 304 & 1.75 & 40.2 & 15.8 & 36.5 & 58.2 \\
Quang Ngai & 303 & 1.08 & 42.4 & 20.2 & 29.4 & 49.8 \\
Phu Tho & 284 & 1.83 & 49.6 & 24.0 & 25.7 & 44.4 \\
Gia Lai & 238 & 1.69 & 49.8 & 25.8 & 24.4 & 45.8 \\
Hau Giang & 211 & 0.99 & 39.3 & 20.6 & 25.6 & 48.8 \\
Lao Cai & 208 & 7.06 & 88.5 & 29.5 & 21.6 & 38.9 \\
An Giang & 202 & 2.26 & 55.5 & 31.8 & 17.8 & 36.1 \\
Dong Thap & 180 & 1.44 & 46.8 & 31.0 & 23.3 & 37.8 \\
Quang Tri & 179 & 1.25 & 46.1 & 26.4 & 21.2 & 42.5 \\
Quang Binh & 165 & 2.58 & 37.0 & 18.0 & 30.9 & 54.5 \\
Vinh Long & 159 & 2.56 & 56.9 & 33.8 & 13.2 & 31.4 \\
Son La & 102 & 3.67 & 74.1 & 42.9 & 15.7 & 24.5 \\
\end{longtable}}


{\scriptsize
\setlength{\tabcolsep}{4pt}
\begin{longtable}{lrrrrrr}
\caption{Calibration and confidence per region. Median/trimmed bias near 0 = the model is centered.}\label{tab:regcal}\\
\toprule
\textbf{Region} & \textbf{Raw Bias\%} & \textbf{Med.Bias\%} & \textbf{Trim.Bias\%} & \textbf{Int.Width\%} & \textbf{In-Band\%} & \textbf{HiConf MedAPE\%} \\
\midrule
\endfirsthead
\multicolumn{7}{c}{\tablename\ \thetable\ -- continued}\\
\toprule
\textbf{Region} & \textbf{Raw Bias\%} & \textbf{Med.Bias\%} & \textbf{Trim.Bias\%} & \textbf{Int.Width\%} & \textbf{In-Band\%} & \textbf{HiConf MedAPE\%} \\
\midrule
\endhead
\midrule\multicolumn{7}{r}{\itshape continued on next page}\\
\endfoot
\bottomrule
\endlastfoot
Da Nang & 5.3 & -0.1 & 0.4 & 69 & 77 & 8.8 \\
Ho Chi Minh & 22.6 & -0.7 & 3.3 & 130 & 76 & 18.6 \\
Hai Phong & 5.8 & -0.7 & -0.2 & 96 & 78 & 9.4 \\
Binh Duong & 13.5 & -0.2 & 1.3 & 92 & 78 & 8.6 \\
Dong Nai & 12.5 & -0.2 & 1.6 & 110 & 77 & 10.3 \\
Khanh Hoa & 7.7 & -0.1 & 0.1 & 95 & 76 & 7.9 \\
Ba Ria Vung Tau & 9.5 & -0.5 & 0.0 & 101 & 77 & 7.7 \\
Long An & 11.9 & -0.0 & 2.6 & 144 & 78 & 12.7 \\
Can Tho & 11.2 & -0.5 & 0.3 & 108 & 77 & 9.2 \\
Ha Noi & 127.7 & -2.0 & -0.6 & 107 & 76 & 18.1 \\
Hue & 8.9 & -0.0 & 0.8 & 114 & 78 & 9.1 \\
Dak Lak & 13.8 & -0.3 & 0.8 & 118 & 75 & 10.8 \\
Lam Dong & 14.2 & -1.0 & 2.2 & 149 & 76 & 14.1 \\
Bac Ninh & 11.8 & -1.1 & -1.1 & 116 & 78 & 11.6 \\
Quang Nam & 9.0 & -0.4 & 1.1 & 162 & 77 & 15.1 \\
Quang Ninh & 21.5 & -0.6 & 0.1 & 139 & 74 & 14.7 \\
Binh Thuan & 12.9 & -0.8 & 1.3 & 196 & 77 & 16.9 \\
Tay Ninh & 11.1 & 1.4 & 4.6 & 258 & 75 & 18.1 \\
Thanh Hoa & 19.3 & 1.5 & 5.0 & 199 & 75 & 12.4 \\
other small (46) & 78.4 & 15.6 & 40.0 & 218 & 70 & 39.6 \\
Hoa Binh & 18.5 & 1.3 & 5.9 & 265 & 75 & 20.0 \\
Nghe An & 19.9 & -1.5 & -0.5 & 118 & 79 & 12.5 \\
Vinh Phuc & 16.0 & -1.6 & -0.2 & 134 & 79 & 12.7 \\
Binh Dinh & 8.3 & -1.5 & -0.6 & 115 & 81 & 11.4 \\
Nam Dinh & 5.0 & 0.2 & 1.2 & 122 & 80 & 10.9 \\
Thai Binh & 6.7 & -1.8 & 0.5 & 122 & 76 & 13.3 \\
Hai Duong & 5.4 & -2.5 & -2.5 & 118 & 77 & 9.1 \\
Binh Phuoc & 20.8 & 5.5 & 12.3 & 403 & 72 & 25.3 \\
Bac Giang & 14.2 & 0.3 & 1.9 & 222 & 73 & 21.5 \\
Tien Giang & 10.2 & -0.9 & 0.5 & 226 & 79 & 16.1 \\
Kien Giang & 28.3 & 1.2 & 7.3 & 232 & 71 & 21.8 \\
Hung Yen & 9.1 & -1.8 & -0.6 & 138 & 78 & 11.9 \\
Phu Yen & 12.5 & -0.8 & 1.1 & 180 & 76 & 19.3 \\
Ha Nam & 14.0 & 1.0 & 1.0 & 137 & 72 & 12.9 \\
Tra Vinh & 10.8 & 2.2 & 3.5 & 254 & 84 & 24.0 \\
Ninh Thuan & 35.1 & 2.4 & 10.1 & 221 & 76 & 17.8 \\
Ben Tre & 19.5 & 6.7 & 10.8 & 318 & 76 & 19.0 \\
Thai Nguyen & 16.8 & -0.9 & -0.1 & 168 & 77 & 11.8 \\
Quang Ngai & 23.8 & 5.5 & 8.1 & 207 & 77 & 16.2 \\
Phu Tho & 23.2 & -0.9 & 4.6 & 201 & 70 & 19.5 \\
Gia Lai & 27.0 & 6.4 & 11.2 & 374 & 71 & 17.7 \\
Hau Giang & 18.9 & 10.9 & 11.2 & 524 & 65 & 28.4 \\
Lao Cai & 56.0 & -3.8 & 9.0 & 163 & 69 & 22.9 \\
An Giang & 29.5 & 8.1 & 13.5 & 259 & 75 & 21.2 \\
Dong Thap & 21.9 & 4.8 & 11.8 & 266 & 71 & 21.6 \\
Quang Tri & 21.9 & 3.8 & 9.1 & 296 & 75 & 18.2 \\
Quang Binh & 9.2 & -4.8 & -2.7 & 153 & 70 & 11.6 \\
Vinh Long & 27.6 & 11.6 & 13.4 & 289 & 71 & 25.0 \\
Son La & 27.5 & -9.7 & 2.7 & 212 & 72 & 33.2 \\
\end{longtable}}


\paragraph{How to read these tables (the three key insights).}
\begin{enumerate}
  \item \textbf{Accuracy is not uniform --- and the report says so.} Well-sampled markets are sharp
        (Da Nang $11.3\pct$ median APE, Khanh Hoa $11.2\pct$, Hai Phong / Binh Duong $\sim 13\pct$),
        while \textbf{Ho Chi Minh is the clear weak spot} ($23.7\pct$) --- a high-variance market the
        blended national average was hiding.
  \item \textbf{Every region is unbiased.} In Table~\ref{tab:regcal} the median and trimmed bias sit
        near $0$ for essentially every province: no market is systematically over- or under-valued.
        The large raw MAPE values are the label-noise tail (Section~\ref{sec:errors}), not a regional
        tilt.
  \item \textbf{Confidence tracks difficulty.} Regions with wider mean intervals are exactly the
        higher-error regions, and the \emph{high-confidence} median APE column is markedly sharper than
        the overall one in every region --- the model's uncertainty signal works locally, not just in
        aggregate.
\end{enumerate}

%======================================================================
\section{Error Cases per Region, and a Focus on the Worst}\label{sec:errors}
%======================================================================
For every province, the worst error cases sit at the top of its CSV file. Ranking provinces by the
share of large-error listings (absolute percentage error $>50\pct$) and comparing the price/m\textsuperscript{2}
of their error listings against their normal listings reveals a single, consistent cause.

\begin{table}[h]
\centering
\begin{tabular}{lrccc}
\toprule
\textbf{Region} & \textbf{n} & \textbf{Large-error rate} & \textbf{Err.\ median ppm2} & \textbf{Normal median ppm2} \\
\midrule
Ho Chi Minh & 14,904 & 21.2\pct & 39.4M & 77.6M \\
Long An     &  6,496 & 17.8\pct & 7.5M  & 18.8M \\
Thanh Hoa   &  1,150 & 22.1\pct & 15.0M & 41.0M \\
Binh Thuan  &  1,518 & 20.8\pct & 12.3M & 27.8M \\
\bottomrule
\end{tabular}
\caption{In every region, error listings have a far \emph{lower} price/m\textsuperscript{2} than
normal listings --- the signature of underpriced / mislabeled listings, not model failure.}
\end{table}

\paragraph{The error is in the label, not the model --- shown for every region.}
Table~\ref{tab:regex} lists, for \emph{each} region, its single worst-error listing next to a
representative normal listing, together with the features that explain the gap (area, floors,
price/m\textsuperscript{2}), the listed price, and the model's prediction. The table is meant to be
read row-pair by row-pair, and every pair tells the same story:

\begin{itemize}
  \item \textbf{The discriminating feature is price/m\textsuperscript{2}.} The worst-error listing
        always has an \emph{implausibly low} price/m\textsuperscript{2} (often $1$--$4$M, versus tens
        of millions for the normal listing in the same region). Area and floors are ordinary --- it is
        the price feature that is broken.
  \item \textbf{The prediction comparison proves the point.} For the worst case the \emph{model's}
        value is far above the tiny listed price (e.g.\ tens or hundreds of billions vs.\ a fraction
        of a billion), because the model reconstructs a sensible value from the normal features. For
        the normal case the model's value sits within a few percent of the listed price. The huge
        ``error\%'' is therefore the distance between a good prediction and a bad label --- not model
        failure.
\end{itemize}

{\scriptsize
\setlength{\tabcolsep}{4pt}
\begin{longtable}{lllrrrrrr}
\caption{Worst-error vs.\ a normal listing for every region, with the features that drive the gap. Error listings consistently show an implausibly low price/m$^2$; the model predicts a sensible value, so the large ``error'' reflects a bad listed price, not a bad prediction.}\label{tab:regex}\\
\toprule
\textbf{Region} & \textbf{Case} & \textbf{Subtype} & \textbf{Area} & \textbf{Flr} & \textbf{ppm2(M)} & \textbf{Listed(B)} & \textbf{Model(B)} & \textbf{Err\%} \\
\midrule
\endfirsthead
\multicolumn{9}{c}{\tablename\ \thetable\ -- continued}\\
\toprule
\textbf{Region} & \textbf{Case} & \textbf{Subtype} & \textbf{Area} & \textbf{Flr} & \textbf{ppm2(M)} & \textbf{Listed(B)} & \textbf{Model(B)} & \textbf{Err\%} \\
\midrule
\endhead
\midrule\multicolumn{9}{r}{\itshape continued on next page}\\
\endfoot
\bottomrule
\endlastfoot
Da Nang & worst & normal & 82 & 4 & 1.2 & 0.10 & 9.81 & 9,706 \\
Da Nang & normal & alley frontage & 186 & 3 & 55.9 & 10.39 & 11.10 & 7 \\
Ho Chi Minh & worst & frontage & 160 & 10 & 1.5 & 0.24 & 127.56 & 53,050 \\
Ho Chi Minh & normal & normal & 56 & -- & 89.3 & 5.00 & 5.44 & 9 \\
Hai Phong & worst & frontage & 68 & 2 & 1.5 & 0.10 & 8.74 & 8,642 \\
Hai Phong & normal & alley frontage & 100 & 5 & 108.0 & 10.80 & 11.56 & 7 \\
Binh Duong & worst & normal & 220 & 5 & 1.0 & 0.23 & 13.25 & 5,659 \\
Binh Duong & normal & normal & 64 & 3 & 52.3 & 3.35 & 3.57 & 7 \\
Dong Nai & worst & frontage & 99 & -- & 1.0 & 0.10 & 7.22 & 7,120 \\
Dong Nai & normal & normal & 80 & 2 & 57.5 & 4.60 & 4.90 & 6 \\
Khanh Hoa & worst & alley frontage & 70 & -- & 1.6 & 0.11 & 6.26 & 5,587 \\
Khanh Hoa & normal & frontage & 520 & -- & 180.8 & 94.00 & 99.30 & 6 \\
Ba Ria Vung Tau & worst & frontage & 117 & -- & 1.4 & 0.16 & 10.19 & 6,269 \\
Ba Ria Vung Tau & normal & alley & 54 & 2 & 72.2 & 3.90 & 4.13 & 6 \\
Long An & worst & townhouse project & 430 & -- & 1.8 & 0.77 & 12.61 & 1,539 \\
Long An & normal & townhouse project & 52 & -- & 28.8 & 1.50 & 1.39 & 7 \\
Can Tho & worst & normal & 80 & 3 & 4.0 & 0.32 & 5.57 & 1,640 \\
Can Tho & normal & villa & 400 & -- & 25.0 & 10.00 & 10.65 & 6 \\
Ha Noi & worst & townhouse project & 290 & 8 & 1.0 & 0.30 & 126.41 & 42,035 \\
Ha Noi & normal & alley frontage & 39 & 4 & 111.5 & 4.35 & 4.71 & 8 \\
Hue & worst & frontage & 84 & 3 & 2.6 & 0.22 & 8.48 & 3,754 \\
Hue & normal & frontage & 110 & 3 & 54.5 & 6.00 & 6.37 & 6 \\
Dak Lak & worst & frontage & 122 & 2 & 1.0 & 0.12 & 10.48 & 8,490 \\
Dak Lak & normal & normal & 88 & 1 & 33.5 & 2.95 & 2.73 & 7 \\
Lam Dong & worst & villa & 500 & -- & 1.3 & 0.65 & 17.88 & 2,655 \\
Lam Dong & normal & frontage & 93 & 5 & 263.4 & 24.50 & 22.65 & 8 \\
Bac Ninh & worst & shophouse & 100 & -- & 1.1 & 0.11 & 9.66 & 8,683 \\
Bac Ninh & normal & villa & 180 & -- & 66.7 & 12.00 & 11.22 & 7 \\
Quang Nam & worst & normal & 127 & -- & 2.8 & 0.35 & 3.19 & 810 \\
Quang Nam & normal & townhouse project & 160 & -- & 15.6 & 2.50 & 2.31 & 8 \\
Quang Ninh & worst & villa & 296 & -- & 1.0 & 0.30 & 35.23 & 11,645 \\
Quang Ninh & normal & normal & 2100 & 27 & 314.3 & 660.00 & 706.55 & 7 \\
Binh Thuan & worst & townhouse project & 75 & -- & 2.7 & 0.20 & 3.37 & 1,587 \\
Binh Thuan & normal & frontage & 100 & -- & 24.0 & 2.40 & 2.57 & 7 \\
Tay Ninh & worst & frontage & 174 & -- & 1.3 & 0.22 & 1.44 & 554 \\
Tay Ninh & normal & normal & 232 & -- & 8.6 & 2.00 & 1.85 & 7 \\
Thanh Hoa & worst & frontage & 35 & 3 & 3.7 & 0.13 & 2.67 & 1,953 \\
Thanh Hoa & normal & frontage & 637 & 4 & 45.5 & 29.00 & 27.09 & 7 \\
Hoa Binh & worst & frontage & 436 & 2 & 1.3 & 0.55 & 4.49 & 716 \\
Hoa Binh & normal & villa & 180 & -- & 9.1 & 1.63 & 1.51 & 8 \\
Nghe An & worst & alley frontage & 79 & 4 & 2.2 & 0.18 & 13.12 & 7,354 \\
Nghe An & normal & townhouse project & 128 & -- & 30.5 & 3.90 & 4.18 & 7 \\
Vinh Phuc & worst & frontage & 78 & 4 & 2.3 & 0.18 & 4.91 & 2,625 \\
Vinh Phuc & normal & frontage & 95 & 3 & 121.1 & 11.50 & 10.66 & 7 \\
Binh Dinh & worst & normal & 350 & -- & 1.4 & 0.50 & 9.24 & 1,748 \\
Binh Dinh & normal & frontage & 42 & 2 & 57.1 & 2.40 & 2.58 & 7 \\
Nam Dinh & worst & alley & 469 & -- & 1.5 & 0.70 & 6.84 & 876 \\
Nam Dinh & normal & alley frontage & 40 & 3 & 39.5 & 1.58 & 1.46 & 8 \\
Thai Binh & worst & villa & 240 & -- & 2.6 & 0.62 & 4.17 & 568 \\
Thai Binh & normal & normal & 40 & 3 & 38.8 & 1.55 & 1.66 & 7 \\
Hai Duong & worst & townhouse project & 200 & -- & 2.9 & 0.58 & 5.81 & 901 \\
Hai Duong & normal & alley frontage & 41 & 3 & 59.8 & 2.45 & 2.62 & 7 \\
Binh Phuoc & worst & frontage & 160 & -- & 2.4 & 0.39 & 2.52 & 547 \\
Binh Phuoc & normal & alley frontage & 300 & -- & 1.2 & 0.36 & 0.39 & 8 \\
Bac Giang & worst & townhouse project & 270 & -- & 3.7 & 1.00 & 7.05 & 605 \\
Bac Giang & normal & townhouse project & 150 & 5 & 100.0 & 15.00 & 13.65 & 9 \\
Tien Giang & worst & alley frontage & 160 & -- & 2.1 & 0.33 & 1.91 & 480 \\
Tien Giang & normal & frontage & 151 & 1 & 27.8 & 4.20 & 3.89 & 7 \\
Kien Giang & worst & villa & 13 & -- & 8.0 & 0.10 & 1.20 & 1,085 \\
Kien Giang & normal & normal & 114 & -- & 17.5 & 2.00 & 2.17 & 8 \\
Hung Yen & worst & normal & 62 & -- & 2.8 & 0.17 & 2.17 & 1,149 \\
Hung Yen & normal & normal & 63 & -- & 55.6 & 3.50 & 3.77 & 8 \\
Phu Yen & worst & townhouse project & 1169 & -- & 2.0 & 2.35 & 14.80 & 530 \\
Phu Yen & normal & frontage & 100 & -- & 41.0 & 4.10 & 3.76 & 8 \\
Ha Nam & worst & townhouse project & 165 & -- & 1.7 & 0.28 & 10.87 & 3,844 \\
Ha Nam & normal & townhouse project & 52 & 5 & 101.0 & 5.25 & 4.88 & 7 \\
Tra Vinh & worst & alley & 485 & -- & 1.0 & 0.48 & 3.84 & 691 \\
Tra Vinh & normal & frontage & 76 & 1 & 59.2 & 4.50 & 4.05 & 10 \\
Ninh Thuan & worst & shophouse & 210 & 4 & 2.3 & 0.48 & 5.66 & 1,086 \\
Ninh Thuan & normal & shophouse & 220 & 4 & 13.6 & 3.00 & 3.21 & 7 \\
Ben Tre & worst & alley & 80 & -- & 3.1 & 0.25 & 1.33 & 430 \\
Ben Tre & normal & normal & 120 & -- & 4.6 & 0.55 & 0.60 & 8 \\
Thai Nguyen & worst & frontage & 131 & 1 & 1.3 & 0.17 & 3.37 & 1,931 \\
Thai Nguyen & normal & normal & 40 & 2 & 50.0 & 2.00 & 1.86 & 7 \\
Quang Ngai & worst & alley frontage & 133 & -- & 1.2 & 0.16 & 1.69 & 959 \\
Quang Ngai & normal & frontage & 105 & 3 & 31.0 & 3.25 & 3.51 & 8 \\
Phu Tho & worst & villa & 197 & -- & 2.5 & 0.49 & 4.07 & 721 \\
Phu Tho & normal & frontage & 700 & -- & 9.1 & 6.40 & 6.97 & 9 \\
Gia Lai & worst & townhouse project & 300 & -- & 1.6 & 0.47 & 3.20 & 584 \\
Gia Lai & normal & frontage & 150 & 2 & 19.0 & 2.85 & 2.59 & 9 \\
Hau Giang & worst & frontage & 683 & -- & 1.2 & 0.80 & 4.09 & 412 \\
Hau Giang & normal & villa & 138 & -- & 12.0 & 1.65 & 1.80 & 9 \\
Lao Cai & worst & townhouse project & 120 & 4 & 1.0 & 0.12 & 4.16 & 3,368 \\
Lao Cai & normal & frontage & 145 & -- & 51.0 & 7.40 & 6.75 & 9 \\
An Giang & worst & townhouse project & 200 & -- & 3.0 & 0.59 & 4.08 & 591 \\
An Giang & normal & frontage & 180 & -- & 34.4 & 6.20 & 5.58 & 10 \\
Dong Thap & worst & normal & 150 & -- & 2.7 & 0.41 & 1.93 & 371 \\
Dong Thap & normal & frontage & 184 & -- & 6.0 & 1.10 & 1.02 & 8 \\
Quang Tri & worst & townhouse project & 325 & -- & 1.3 & 0.44 & 3.24 & 641 \\
Quang Tri & normal & normal & 142 & -- & 5.3 & 0.76 & 0.68 & 10 \\
Quang Binh & worst & townhouse project & 200 & -- & 5.0 & 1.00 & 4.97 & 397 \\
Quang Binh & normal & townhouse project & 50 & -- & 42.0 & 2.10 & 2.28 & 9 \\
Vinh Long & worst & frontage & 361 & -- & 1.0 & 0.37 & 4.11 & 1,010 \\
Vinh Long & normal & shophouse & 44 & -- & 15.2 & 0.66 & 0.74 & 12 \\
Son La & worst & villa & 190 & -- & 1.9 & 0.37 & 3.00 & 710 \\
Son La & normal & townhouse project & 85 & 4 & 51.8 & 4.40 & 4.76 & 8 \\
\end{longtable}}


\subsection*{Focus: Ho Chi Minh City --- the region with the most errors}
Ho Chi Minh is both the largest major market (14,904 test listings) and the one producing the most
error cases ($21.2\pct$ large-error rate, $\approx 3{,}160$ listings). Its defensible interpretation:

\begin{itemize}
  \item \textbf{The worst case is a data error, and the model caught it.} A 160 m\textsuperscript{2}
        frontage house listed at $1.5$M/m\textsuperscript{2} ($0.24$B total) is impossible in Ho Chi
        Minh, where normal frontage houses trade near $80$--$100$M/m\textsuperscript{2}. The model
        valued it at $\sim 128$B --- almost certainly the true order of magnitude --- producing a
        $53{,}000\pct$ ``error'' that is entirely the fault of a mistyped price. This listing is
        exactly what an AVM should \emph{flag}, not what it should be judged for missing.
  \item \textbf{Error listings are systematically cheap.} Their median price/m\textsuperscript{2} is
        $39.4$M versus $77.6$M for normal listings --- error cases are dominated by implausibly
        under-priced posts, and $19.9\pct$ sit in the extreme-price tail versus $8.8\pct$ of normal
        listings.
  \item \textbf{Genuine market variance, honestly reported.} Beyond bad labels, Ho Chi Minh is a
        high-dispersion metro: prices for superficially similar listings vary widely by exact
        street and alley. This is why its intervals are wider and more of its listings route to
        Tier~C. The right response is district-level modeling, and the segmentation makes that need
        visible instead of averaging it away.
  \item \textbf{The product remains safe here.} On the higher-confidence half of Ho Chi Minh
        listings the typical error is $\sim 18\pct$, and the anomaly checks (interval width, asking
        price outside band) route the problematic listings to review. Normal Ho Chi Minh listings ---
        such as the $56$ m\textsuperscript{2} home valued to $8.8\pct$ above --- are served sharply.
\end{itemize}

%======================================================================
\section{Limitations and Next Steps}
%======================================================================
\begin{itemize}
  \item \textbf{Thin regions stay uncertain.} Small provinces have wider intervals and larger errors;
        the policy correctly routes more of them to Tier~C. More regional data is the only real fix.
  \item \textbf{The error tail is contained, not cleaned.} Flagging limits the damage; cleaning the
        underlying suspicious-price listings is the upstream fix and the single biggest lever on raw
        MAPE.
  \item \textbf{Random split.} All numbers use a random test split; a time-based split (hold out the
        most recent months) is the next step to prove the model holds on \emph{future} listings.
  \item \textbf{Ho Chi Minh} is the clearest modeling follow-up --- a high-variance metropolitan
        market that likely needs district-level segmentation.
\end{itemize}

\vspace{1em}
\hrule
\vspace{0.5em}
\noindent\footnotesize All figures are computed from the held-out test predictions of the production
(shrinkage-blend) model. Per-region data lives in \texttt{Rule-based Scene/Evaluation/regional\_report/regions/};
the full machine-readable report is regenerated with \texttt{python3 src/build\_regional\_report.py}.

\end{document}
